% *======================================================================*
%  Cactus Thorn template for ThornGuide documentation
%  Author: Ian Kelley
%  Date: Sun Jun 02, 2002
%  $Header: /cactusdevcvs/Cactus/doc/ThornGuide/template.tex,v 1.12 2004/01/07 20:12:39 rideout Exp $
%
%  Thorn documentation in the latex file doc/documentation.tex
%  will be included in ThornGuides built with the Cactus make system.
%  The scripts employed by the make system automatically include
%  pages about variables, parameters and scheduling parsed from the
%  relevant thorn CCL files.
%
%  This template contains guidelines which help to assure that your
%  documentation will be correctly added to ThornGuides. More
%  information is available in the Cactus UsersGuide.
%
%  Guidelines:
%   - Do not change anything before the line
%       % START CACTUS THORNGUIDE",
%     except for filling in the title, author, date, etc. fields.
%        - Each of these fields should only be on ONE line.
%        - Author names should be separated with a \\ or a comma.
%   - You can define your own macros, but they must appear after
%     the START CACTUS THORNGUIDE line, and must not redefine standard
%     latex commands.
%   - To avoid name clashes with other thorns, 'labels', 'citations',
%     'references', and 'image' names should conform to the following
%     convention:
%       ARRANGEMENT_THORN_LABEL
%     For example, an image wave.eps in the arrangement CactusWave and
%     thorn WaveToyC should be renamed to CactusWave_WaveToyC_wave.eps
%   - Graphics should only be included using the graphicx package.
%     More specifically, with the "\includegraphics" command.  Do
%     not specify any graphic file extensions in your .tex file. This
%     will allow us to create a PDF version of the ThornGuide
%     via pdflatex.
%   - References should be included with the latex "\bibitem" command.
%   - Use \begin{abstract}...\end{abstract} instead of \abstract{...}
%   - Do not use \appendix, instead include any appendices you need as
%     standard sections.
%   - For the benefit of our Perl scripts, and for future extensions,
%     please use simple latex.
%
% *======================================================================*
%
% Example of including a graphic image:
%    \begin{figure}[ht]
% 	\begin{center}
%    	   \includegraphics[width=6cm]{MyArrangement_MyThorn_MyFigure}
% 	\end{center}
% 	\caption{Illustration of this and that}
% 	\label{MyArrangement_MyThorn_MyLabel}
%    \end{figure}
%
% Example of using a label:
%   \label{MyArrangement_MyThorn_MyLabel}
%
% Example of a citation:
%    \cite{MyArrangement_MyThorn_Author99}
%
% Example of including a reference
%   \bibitem{MyArrangement_MyThorn_Author99}
%   {J. Author, {\em The Title of the Book, Journal, or periodical}, 1 (1999),
%   1--16. {\tt http://www.nowhere.com/}}
%
% *======================================================================*

% If you are using CVS use this line to give version information
% $Header: /cactusdevcvs/Cactus/doc/ThornGuide/template.tex,v 1.12 2004/01/07 20:12:39 rideout Exp $

\documentclass{article}

% Use the Cactus ThornGuide style file
% (Automatically used from Cactus distribution, if you have a
%  thorn without the Cactus Flesh download this from the Cactus
%  homepage at www.cactuscode.org)
\usepackage{../../../../doc/latex/cactus}
\usepackage{color}

\begin{document}

% The author of the documentation
\author{Roland Haas \textless roland.haas@physics.gatech.edu\textgreater}

% The title of the document (not necessarily the name of the Thorn)
\title{RiemannSolverHLL}

% the date your document was last changed, if your document is in CVS,
% please use:
%    \date{$ $Date: 2004/01/07 20:12:39 $ $}
\date{April 08 2010}

\maketitle

% Do not delete next line
% START CACTUS THORNGUIDE

% Add all definitions used in this documentation here
%   \def\mydef etc

% Add an abstract for this thorn's documentation
\begin{abstract}
    This thorn solves the Riemann problem in magnetohydrodynamics using a slab
    geometry approach. It emplys the HLL scheme by Harten, Lax and
    van~Leer~\cite{hll}.
    It currently computes the required eigenvalues using the quadratic
    equation approximation used in the ECHO code~\cite{echo}.
\end{abstract}

% The following sections are suggestive only.
% Remove them or add your own.

\section{Introduction}
If this thorn ever grows up I promise I will write some rivetting indroduction
to why Riemann solvers are needed. For now I'll state for the record that this
is the simplest possible Riemann solver for MHD and seems to work fine as it
is used by all groups (ie. Whisky, the group at UIUC) even though least the
first group has access to a full Riemann solver and a Roe like solve could
also be written. Personal communication by Bruno Giacommazzo tells me that the
important thing is the reconstruction, not the Riemann solver. If we want to
go to a better solver I would suggest going to HLLD
anyway~\cite{jets_and_stars}.

\section{Physical System}
For an in-depth description of Riemann solvers see the book by
Toro~\cite{toro} (who seems to be the authority in the field) or his chapter
in~\cite{jets_and_stars}.

The Riemann problem in 1D hydrodynamics (and we reduce everything to 1D by
using a slab geometry approach) starts with initial data such that the initial
conditions are constant (but different) on each side of an interface. In terms
of the prototypical advection equation that seems to be used universally for
these things it reads
\begin{gather}
    \partial_t u + a \partial_x u = 0 \\
    u(x,0) = u_0(x) = 
      \begin{cases}
          % someone decided that defining \def\text#1{{\rm #1}} is a very good
          % idea. This interferes with AMSMaths \text command meaning I can
          % not use it here. 
          u_L & \mbox{if $x < 0$}, \\
          u_R & \mbox{if $x > 0$},
      \end{cases}
\end{gather}
as shown below.
\begin{center}
\input{riemann_problem.tex}
\end{center}

For the case of the advection equation shown the solution is easy to find and
is
\begin{equation}
    u(x,t) = u_0(x-a t) = \begin{cases} 
                          u_L & \mbox{if $x - at < 0$} \\
                          u_R & \mbox{if  $x -at > 0$}
                          \end{cases}.
\end{equation}
And even for the classical hydrodynamics equations written in conservative
form, the solution is known analytically. 

For magnetohydrodynamics things are more complicated and no analytical
solution is known (however Bruno Giacomazzo implement a full solver for MHD
numerically).

The equivalent of the advection equation in MHD reads (from~\cite{whiskymhd})
\begin{equation}
    \frac{1}{\sqrt{-g}} \left[ \partial_t (\sqrt{\gamma} \vec F^0) +
    \partial_i (\sqrt{-g} F^i)\right] = \vec S
\end{equation},
for some term $\vec S$ and state vector $F^i = [D, S_j, \tau, B^k]^T$ and
fluxes
\begin{equation}
    \vec F^i = \left( \begin{gathered}
        D \tilde v^i / \alpha \\
        S_j \tilde v^i + (p + b^2/2) \delta^i_j - b_j B^i/W \\
        \tau \tilde v^i / \alpha + (P + b^2/2) v^i - \alpha b^0 B^i / W \\
        B^k \tilde v^i / \alpha - B^i \tilde v^k / \alpha
    \end{gathered} \right).
\end{equation}
The quantities are all defined in~\cite{whiskymhd} and are as used in the
code:
\begin{align}
    D &= \rho W, \\
    S_j &= (\rho h + B^2 ) W^2 v_j - \alpha b^0 b_j, \\
    \tau &= (\rho h + b^2) W^2 - (P+b62/2) - (\alpha b^0)^2 - D, \\
    B^i &= {}^*F^{i \nu} n_\nu.
\end{align}
A useful set of equations links the magnetic field as measured by the
observers comiving with the fluid $b^\mu$ to that measured by the Eulerian
observers $B^i$
\begin{equation}
    b^0 = \frac{W B^i v_i}{\alpha}, \quad b^i = \frac{B^i + \alpha b^0
    u^i}{W}, \quad b^2 = b_\mu b^\mu = \frac{B^2 + (\alpha b^0)^2}{W^2}.
\end{equation}
as well as the fact that $W v_i = u_i$.

The HLL solver computes numerical fluxes according to 
\begin{equation}
    \vec F^i = \frac{\lambda_R \vec F^i_L - \lambda_L \vec F^i_R +
    \lambda_R*\lambda_L (F^0_R-F^0_L)}{\lambda_R - \lambda_L}
\end{equation}
where $\lambda_L$ and $\lambda_R$ are the fastest left and right-going speeds
as defined in~\cite{echo}. Note that the sign convention used differs from
that of~\cite{whiskymhd}. Above $\vec F^i_L$ and $\vec F^i_R$ are the fluxes
to the left and right of the interface respecively and 
\begin{eqnarray}
    \lambda_L = min(0,\lambda_{L,left}, \lambda_{L,right}),
    \lambda_R = max(0,\lambda_{R,left}, \lambda_{R,right}).
\end{eqnarray}

We compute the wave speeds using the quadratic approximation of~\cite{echo} as
\begin{equation}
    \lambda^{\prime x}_{L/R} = \frac{(1-a^2) v^x \mp 
                       \sqrt{a^2 (1-v^2) [(1-v^2a^2) \gamma^{xx} -
                                          (1-a^2)(v^x)^2]}
                      }{1-v^2a^2},
\end{equation}
where 
\begin{equation}
    a^2 = c_s^2 + c_a^2 - c_s^2 c_a^2,
\end{equation}
and
\begin{equation}
    c_s^2 = (\rho\left(\frac{dP}{d\rho}\right)_\epsilon + P
    \left(\frac{dP}{d\epsilon}\right)/rho) / (\rho \, h),
    \qquad
    c_a^2 = \frac{b^2}{\rho\, h + b^2},
\end{equation}
are the speed of sound and Alf\'en speed respectively. Finally we transform
from these Lagrangian speeds to Eulerian speeds via
\begin{equation}
    \lambda_{L/R} = \alpha \lambda^{\prime x}_{L/R} - \beta^x.
\end{equation}

\section{Numerical Implementation}
Not much to say here. We compute signal speeds in \verb|signal_speeds.c|,
the individual fluxes in \verb|flux.c| and combine them according to the HLL
prescription in \verb|hll.c|. The only tricky part is that we only implement
the solver for the $x$ direction and use it for all three flux directions by
permuting the vector and tensor indices (using pointers) and \verb|Whisky's|
\verb|xoffset| etc. variables to offset in the flux direction for the
intercell fluxes.

\section{Using This Thorn}
Activate in ActiveThorns, set \verb|Whisky::riemann_solver| to ``hll''.

\subsection{Obtaining This Thorn}
Currently private to GT. 

\subsection{Basic Usage}
Use like any other Riemann thorn.

\subsection{Special Behaviour}
None that I would know of. Will collect all featurized bugs in here later.

\subsection{Interaction With Other Thorns}
Obviously requires Whisky.

\subsection{Examples}
You bet.

\subsection{Support and Feedback}
Yell at me if it does not work.

\section{History}
None yet.

\subsection{Acknowledgements}
My Mom, my Dad, Jesus, my pet guinea pig, the pope. Sheesh, does this file
ever end?


\begin{thebibliography}{9}
\bibitem{hll} A.~Harten, P.D.~Lax and B.~van~Leer,
  ``On Upstream Differencing and Godunov-Type Schemes for Hyperbolic
  Conservation Laws,''
  SIAM Review {\bf 25}, 35--61 (1983)
  [\url{http://dx.doi.org/10.1137/1025002}]
\bibitem{whiskymhd} B.~Giacomazzo and L.~Rezzolla,
  ``WhiskyMHD: a new numerical code for general relativistic
  magnetohydrodynamics,''
  Class.\ Quant.\ Grav.\  {\bf 24}, S235 (2007)
  [arXiv:gr-qc/0701109].
  %%CITATION = CQGRD,24,S235;%%
\bibitem{echo} L.~Del Zanna, O.~Zanotti, N.~Bucciantini and P.~Londrillo,
  ``ECHO: an Eulerian Conservative High Order scheme for general relativistic
  magnetohydrodynamics and magnetodynamics,''
  arXiv:0704.3206 [astro-ph].
  %%CITATION = ARXIV:0704.3206;%%
\bibitem{toro} E.~F.~Toro, 
    ``Riemann solvers and numerical methods for fluid dynamics: a practical
    introduction,''
    2nd ed. Berlin: Springer, 1999. 
\bibitem{jets_and_stars} S.~Massaglia (ed.),
    ``Jets from young stars III: numerical MHD and instabilities,''
    Berlin: Springer, 2008. 




\end{thebibliography}

% Do not delete next line
% END CACTUS THORNGUIDE

\end{document}
